\documentclass[10pt,a4paper]{amsart}
\usepackage[margin=19mm]{geometry}
\usepackage{amsmath,amssymb,mathtools}
\usepackage[scr=rsfs,cal=euler]{mathalfa}
\usepackage{xfrac}
\usepackage[unicode,hidelinks,bookmarks=false]{hyperref}
\newcommand{\ie}{i.e.\ }
\newcommand{\cf}{cf.\ }
\newcommand{\resp}{respectively\ }
\newcommand{\Subsection}{\subsection}
\providecommand{\nobreakdash}{\nobreak\hbox{-}}
\setlength{\emergencystretch}{2em}
\multlinegap0pt
\usepackage{bm}
\usepackage{bbm}
\usepackage{dsfont}
\usepackage{xspace}
\usepackage{cancel}
\usepackage{mathtools}
\usepackage{amsthm}
\makeatletter
\@ifundefined{theorem}{\newtheorem{theorem}{Theorem}}{}
\@ifundefined{lemma}{\newtheorem{lemma}[theorem]{Lemma}}{}
\@ifundefined{proposition}{\newtheorem{proposition}[theorem]{Proposition}}{}
\makeatother
\makeatletter
\newcommand{\RR}{\mathbb{R}}
\newcommand{\cB}{\mathcal{B}}
\newcommand{\cD}{\mathcal{D}}
\newcommand{\cN}{\mathcal{N}}
\@ifundefined{proofof}{\newenvironment{proofof}[1][{\hspace{-\blank}}]{{\medskip\noindent\textit{Proof~{#1}}.\ }}{\qed}}{}
\makeatother
\title[Optimal Codes for Deterministic Identification over Gaussian]{Optimal Codes for Deterministic Identification over Gaussian Channels: Closing the Capacity Gap}
\author{Pau Colomer, Christian Deppe, Holger Boche, Andreas Winter}
\date{Source: arXiv:2604.11782v1 (CC BY 4.0)}
\begin{document}
\maketitle

\noindent\emph{Source and licence.} Optimal Codes for Deterministic Identification over Gaussian Channels: Closing the Capacity Gap. Version arXiv:2604.11782v1 (CC BY 4.0), distributed under the Creative Commons Attribution 4.0 International licence, as recorded in the arXiv licence metadata for this exact version and shown on the abstract page https://arxiv.org/abs/2604.11782v1. Full rights notice: licenses/arxiv-2604.11782-CC-BY-4.0.txt.\par\smallskip

\noindent\textit{Editorial scope.} Counted unit: the Theorem whose source label is \texttt{thm:RR} in the article (source0.tex lines 895--900, source label \texttt{thm:RR}), delivered together with the article's own complete original proof of it (source0.tex lines 901--962, one \texttt{proof} environment of 62 lines). No mathematics is written, repaired or re-derived: statement and proof are the article's own text line for line. Required context retained uncounted: prop:projection (lines 302--307); lemma:angle (lines 340--345); prop:size\_bound (lines 381--386); eq:encoding\_code\_size (lines 682--684); eq:decoder (lines 695--697); eq:layer\_error1 (lines 731--736); eq:main\_Pe2 (lines 758--760). Every definition, display and label of those lines is retained unchanged. Omitted from this package and not redistributed: 1--301, 308--339, 346--380, 387--681, 685--694, 698--730, 737--757, 761--894, 963--1119. Nothing inside a retained excerpt is omitted or altered: every retained body is delivered exactly as the article's lines are written, so no omission receipt of any kind accompanies this package. Citations: the retained excerpts carry no citation command at all, so no bibliography entry of the article is transcribed and no reference is redistributed. Reference adaptation: none was needed and none was made; every reference inside the delivered unit resolves inside it, so no unresolved reference or citation is produced. Printed slips kept literal: no printed word, symbol, hypothesis or number of the counted statement or of its proof was altered. Layout and class: the article's own class is \texttt{IEEEtran} with packages including braket, amsmath, amssymb, amsfonts, dsfont, amsthm, algorithmic, graphicx, textcomp, xcolor, comment, cite, flushend; the wrapper is the lane's pinned \texttt{amsart} editorial wrapper at 10pt on A4, which restates the theorem-like environments the retained text needs and adds the article's own macro definitions that the retained text needs (\texttt{\textbackslash RR}, \texttt{\textbackslash cB}, \texttt{\textbackslash cD}, \texttt{\textbackslash cN}), with extra wrapper packages (bm, bbm, dsfont, xspace, cancel, mathtools). The delivered numbering is the unit's own. Assets: no retained span contains a figure environment, a graphics-inclusion command, a PDF-page inclusion, a table or a TikZ picture; no image file is copied into this package, the package declares no asset dependency and no image file is redistributed with it. Licence: version arXiv:2604.11782v1 of this article is distributed under the Creative Commons Attribution 4.0 International licence, as recorded in the arXiv licence metadata for this exact version and shown on the abstract page https://arxiv.org/abs/2604.11782v1. A full rights notice is delivered at licenses/arxiv-2604.11782-CC-BY-4.0.txt. Characters: every retained excerpt is pure ASCII, so the delivered engine, XeLaTeX with its default Latin Modern text font, reports no missing character.
\begin{proposition}\label{prop:projection}
    Let $\vec{u}\in\RR^n$ be a vector in a high dimensional space and $\vec{Z}=(Z_1,\dots,Z_n)$ a random noise vector where $Z_1,\dots,Z_n\sim\cN(0,1)$ are i.i.d. standard normal random variables. The projection of the noise vector onto  $\vec{u}$ satisfies:
    \[
    \Pr\left(\|\Pi_{\vec{u}}\vec{Z}\|\geq x\right)=2\Phi(-x)\leq\frac{1}{\sqrt{2\pi}}\frac{1}{x}e^{-x^2/2}.
    \]
\end{proposition}

\begin{lemma}\label{lemma:angle}
The minimum separation angle $\theta$ between any two different points $u_j\neq u_k$ in a $d$-angle-dense arrangement on the surface of the $r$-radius ball $\cB_{\vec o}(r,n)$ satisfies
\[
\sin\frac{\theta}{2} \geq \sqrt{\frac{d}{2r}} .
\]
\end{lemma}

\begin{proposition}\label{prop:size_bound}
The number $N$ of elements in a $d$-angle-dense arrangement on a ball $\cB_{\vec o}(r,n)$ of radius $r$ and dimension $n$ satisfies
\[
\log N \geq \frac{n}{2}\log\left(\frac{2r}{d}\right).
\]
\end{proposition}

\begin{equation}\label{eq:encoding_code_size}
\log N_\ell\geq\frac{n+1-\ell}{2}\log \left(\frac{2r_\ell}{d}\right),
\end{equation}

\begin{equation}\label{eq:decoder}
\cD_{i}^{(\ell)}=\{\vec{y}\in\RR^n:\|\Pi_{s_{\ell-1}\to s_{\ell}}\vec{y}-\vec{o}_{s^\ell}\|\leq\sigma\log n\},
\end{equation}

\begin{align}
\!\!\Pr&(\cD_i^{(\ell)}|\vec{u}_i)=\Pr\left(\|\Pi_{s_{\ell-1}\to s_{\ell}}\vec{Y}-\vec{o}_{s^\ell}\|\leq\sigma\log n|\vec{o}_{s^{L}}\right)\nonumber \\
&=\Pr\left(\|\Pi_{s_{\ell-1}\to s_{\ell}}\vec{o}_{s^{L}}-\vec{o}_{s^\ell}+\sigma \Pi_{s_{\ell-1}\to s_{\ell}}\vec{Z}\|\leq \sigma\log n\right)\nonumber\\
&=\Pr\left(\|\Pi_{s_{\ell-1}\to s_{\ell}}\vec{Z}\|\leq\log n\right)\label{eq:layer_error1}\\
&=1-2\Phi(-\log n)=1-\lambda,\nonumber
\end{align}

\begin{equation}\label{eq:main_Pe2}
P_{e,2}(\vec{u}_i|\vec{u}_j)=\Pr\left(\cD_i|\vec{u}_j\right)\leq\Pr\left(\cD_i^{(\ell)}|\vec{u}_j\right)\leq\lambda:=\lambda_2.
\end{equation}

\begin{theorem}\label{thm:RR}
Given an AWGN channel with input power constraint $P>0$ and noise variance $\sigma^2>0$, one can construct an $(n,N,L2^{-nE(n)},2^{-nE(n)})$ DI code for all $E(n)\leq E_0:=\frac{9P}{\sigma^2}$ with $L\gg1$ angle-dense layers achieving a linear rate of
\[
R\geq \frac12\log\frac{1}{E(n)}+\frac12\log\frac{P}{9L\sigma^2}.
\]
\end{theorem}

\begin{proof}
Start by observing that, in our construction, the error performance is dictated by the choice of the accepted range in the decoders at each layer, let us call it $\sigma x(n)$ [see Equation~\eqref{eq:decoder}, there $x(n)=\log n$]. Then, the first type of error at each layer is bounded by $\lambda=2\Phi[-x(n)]$, cf.~Equation~\eqref{eq:layer_error1}. Recalling Proposition~\ref{prop:projection}, and by choosing $x(n)=\sqrt{2nE(n)}$ the following exponential error expression is obtained:
\begin{equation}
\lambda=2\Phi\left(-\sqrt{2nE(n)}\right)\leq 2^{-nE(n)}.
\end{equation}
The total probability of missed identification in a code of $L$ layers is therefore bounded by $\lambda_1\leq L\lambda$. Then, for any arbitrarily large constant $L$, the error converges to 0 when $n\to\infty$.

The second type of error (false identification) can be similarly bounded. Choosing $d=3\sigma x(n)$ one observes [cf.~Equation~\ref{eq:main_Pe2}], that $\lambda_2:=\lambda=2^{-nE(n)}$. Thus, the choices of $x(n)=\sqrt{2nE(n)}$ and $d=3\sigma x(n)$ produce a code with errors that can be written in the desired exponential form. The decoder at each layer is:
\begin{equation}\label{eq:RR_decoder}
\cD_{i}^{(\ell)}=\left\{\vec{y}\in\RR^n:\|\Pi_{s_{\ell-1}\to s_{\ell}}\vec{y}-\vec{o}_{s^\ell}\|\leq\sigma\sqrt{2nE(n)}\right\}.
\end{equation}

Let us study now how the previous definitions fix the code structure. From Lemma~\ref{lemma:angle} one can extract 
\begin{equation}\label{eq:sin1}
    \sin \theta_\ell\simeq\sqrt{\frac{2d}{r_\ell}},
\end{equation}
where we have used the small-angle approximation $\sin\frac{\theta_{\ell}}{2}\approx\frac12\sin\theta_{\ell}$, which becomes asymptotically exact as $n\to\infty$ since $d=o(r_\ell)$, i.e.~$r_\ell$ grows strictly faster than $d$ with the block length. 

Similarly, to make the second type of error work with the previous choice of $d$, it is necessary that $\delta\leq \sigma x(n)$, such that $\|\Pi_{s_{\ell-1}\to s_{\ell}}\vec{o}_{\tilde{s}^{L}}-\vec{o}_{s^\ell}\|\geq d-\delta\geq 2\sigma\log n$. This is ensured if 
\begin{equation}\label{eq:sin2}
\sqrt{L}r_{\ell+1}\sin\theta_\ell=\sigma x(n)
\end{equation}
Then, equating the expressions \eqref{eq:sin1} and \eqref{eq:sin2}, one obtains the following relation between the radius at one layer and the radius at the next:
\begin{equation}\label{eq:RR_hierarchy}
    r_{\ell+1}=\sqrt{\frac{r_\ell\sigma x(n)}{6L}}.
\end{equation}
If one now fixes $r_1=\sqrt{\frac{Pn}{2}}$ as the radius at the first layer, the expression above defines the radius at all other layers. The following relation is obtained:
\begin{equation}\label{eq:RR_radius}
r_\ell = \sqrt{\rho n E(n)}\left[\frac{P}{2\rho E(n)}\right]^\frac{1}{2^\ell}\!\!,\,\,\,\,\text{with}\,\,\,\rho:=\left(\frac{\sigma}{3\sqrt{2}L}\right)^2\!\!.
\end{equation}

As the words have to be inside the input space, it is required that $\sum_{\ell=1}^L r_\ell^2\leq nP$, which is necessary fulfilled if $L r_2^2\leq nP-r_1^2=\frac{nP}{2}$. Then, as
\begin{equation}
L r_2^2=L\rho n E(n)\sqrt{\frac{P}{2\rho E(n)}}=n\frac{\sigma}{6}\sqrt{PE(n)},
\end{equation}
it is necessary that
\begin{equation}
\frac{\sigma}{6}\sqrt{PE(n)}\leq\frac{P}{2}\implies E(n)\leq E_0:=\frac{9P}{\sigma^2}.
\end{equation}

One can use the considerations above to create a good DI code. Let us now analyse its rate performance as a function of the error exponent $E(n)$. Similarly as in the previous section, the rate can be bounded using Equation~\eqref{eq:encoding_code_size} (extracted from Proposition~\ref{prop:size_bound}) and the radius at each layer, given in Equation~\eqref{eq:RR_radius}. Then,
\begin{equation}
\log N_\ell\geq\frac{n-(\ell-1)}{2}\log\left[\frac{2\sqrt{\rho n E(n)}}{3\sigma\sqrt{2nE(n)}}\left(\frac{P}{2\rho E(n)}\right)^\frac{1}{2^\ell}\right].
\end{equation}
Separate now the right-hand side of the equation above into a part with the linear prefactor and a part without it $\frac{n}{2}\log(\dots)-\frac{\ell-1}{2}\log(\dots)$. Then, the negative (second) term can be bounded by the error exponent lower bound $E(n)\geq\frac1n$ and the expression becomes:
\begin{equation}
\begin{split}
\log N_\ell&\geq\!\frac{n}{2}\left[\frac{1}{2^\ell}\log\left(\frac{9PL^2}{\sigma^2 E(n)}\right)\!-\!\log\left(9L\right)\right]\!-\!O\left[(\log n)^\frac{1}{2^\ell}\right]\!.
\end{split}
\end{equation}
The linear rate will then be bounded by:
\begin{align}
R(n):=\sum_{\ell=1}^L\frac{\log N_\ell}{n}\geq\, &\frac12\left(\sum_{\ell=1}^L\frac{1}{2^\ell}\right)\!\!\left(\log\frac{1}{E(n)}+\log\frac{P}{\sigma^2}\right)\nonumber\\
&-\frac{\log 9L}{2}-O\left(\frac{\sqrt{\log n}}{n}\right).\label{eq:RR_rate}
\end{align}

Now, for any $\epsilon>0$ there exists a threshold value $n_0$ and large enough $L$ such that for all $n>n_0$ one finds 
\begin{equation}
R(n)\geq\frac12 \left(\log\frac{1}{E(n)}+\log\frac{P}{9L\sigma^2} \right)-\epsilon.
\end{equation}
Therefore, the claimed rate is indeed achievable.
\end{proof}

\end{document}
